% DES-LOC Appendix - Reconstructed from PDF
% Claude-1 (M001-M025): Key appendix sections
% Appendices A, C, F, G, H, I, J, K, L, M
% (Appendices B, D, E contain extensive figures/proofs → deferred to Claude 9-14)

\appendix

%=============================================================================
\section{Experimental Details and Optimizer Hyperparameter Sweeps}
\label{sec:appendix_a}
%=============================================================================

Here we provide additional experimental details complementing those in Section~4.1, including: a) model architecture details and hyperparameters independent of optimizer choice (Section~\ref{sec:arch_details}), b) our hyperparameter sweep procedure to select optimizer-specific settings (Section~\ref{sec:sweep}), and c) the optimal hyperparameters with those used in Section~5 highlighted in bold.

\subsection{Architecture Details and Hyperparameters}
\label{sec:arch_details}

\begin{table}[h]
\centering
\caption{Model architecture and training parameters. We denote the number of transformer blocks by \#Blocks, number of attention heads by \#Heads, embedding dimension by $d_{\text{model}}$, vocabulary size by $|V|$, and feedforward-layer expansion by Exp.\ Ratio. All models use positional embeddings~\citep{su2024roformer}, the silu activation function, and norm-based gradient clipping with clip-bound $\rho$. Global batch size (summed across all workers) is $|B_G|$, and sequence length is standard for models at these scales. For model initialization we use $\sigma = 1/\sqrt{d_{\text{model}}}$. The total number of steps is denoted by $T$.}
\label{tab:arch}
\begin{tabular}{lccccccccccc}
\toprule
Model Size & Blocks & $d_{\text{model}}$ & $|V|$ & \#Heads & Exp.\ Ratio & ROPE $\theta$ & ACT & Init $\sigma$ & $\rho$ & Seq Len & $|B_G|$ & $T$ \\
\midrule
135M & 30 & 576  & 50K & 9  & 4 & 10000 & silu & 0.04 & 1.0 & 2048 & 1024 & 1536, 3072 \\
720M & 12 & 2048 & 50K & 16 & 4 & 10000 & SiLU & 0.02 & 1.0 & 2048 & 512  & 38912 \\
1.7B & 24 & 2048 & 50K & 16 & 4 & 10000 & silu & 0.02 & 1.0 & 2048 & 1024 & 20480 \\
\bottomrule
\end{tabular}
\end{table}

Table~\ref{tab:arch} summarizes the architectural details of our models, following established practices for large
language models at their respective scales. Unless otherwise stated, we adopt the hyperparameters
recommended by~\citet{allal2025smollm2} for both the 135M and the 1.7B models. We operate at a
batch size of 2M tokens, which is very large for the 135M model at the length of training we
perform~\citep{zhang2025critical} and industry-standard for the 1.7B model~\citep{touvron2023llama2}, we
chose to operate at large batch sizes because adaptive optimizers provide benefits primarily in large-batch training regimes~\citep{kunstner2023noise}. Moreover, we intend \texttt{DES-LOC} for use in cross data-center scenarios, where effectively utilizing available accelerators naturally demands large batch sizes and/or model scales. For both model sizes, we train for approximately $2\times$ the compute-optimal token budget~\citep{hoffmann2022training}, placing our evaluations within the context of extended-duration foundation model training~\citep{allal2025smollm2}. Our chosen token budget is conservative due to resource constraints; for comparison, \citet{allal2025smollm2} used 11 trillion tokens which is over $4000\times$ compute-optimal for the 135M model, and $300\times$ for the 1.7B.

We select warmup and decay schedules following recommendations from~\citet{zhang2025critical,hagele2024scaling,allal2025smollm2}. For the 135M model, the warmup period is set to $T_{\text{WARM}} = 512$
steps, corresponding to roughly 40\% of the compute-optimal training tokens recommended by~\citet{zhang2025critical}. For the 1.7B model, we use the recommended $T_{\text{WARM}} = 2048$ steps from~\citet{allal2025smollm2}, roughly 10\% of total training. The stable-decay period uses a $1 - \text{SQRT}$ schedule over
the final $T_{\text{DECAY}} = 10\% \times T$ steps~\citep{hagele2024scaling}. For shorter runs, such as $T = 1536$ during heterogeneous-data evaluations, we keep the warmup fixed and proportionally scale the decay to ensure well-conditioned parameter updates during the stable learning rate period. The seeds we use for data sampling and for controlling the training algorithms and model are provided in the code accompanying the appendix.

\subsection{Optimizer Parameters Sweeping Procedure}
\label{sec:sweep}

As detailed in Section~\ref{sec:desloc} and verified empirically in Section~\ref{sec:rq2}, the choice of decay rates $\beta_1, \beta_2$
strongly influences the effective synchronization frequencies achievable by both DES-LOC and
Local Adam. This relationship arises directly from the half-life of optimizer states, given by
$\tau_{0.5} = \frac{\ln(0.5)}{\ln(\beta)}$.

For Adam, prior studies such as \citet{wortsman2023stable} have demonstrated a critical interplay between the learning rate ($\eta$), batch size, and the second-momentum decay $\beta_2$. Specifically, increasing either the learning rate or batch size typically demands a lower $\beta_2$ to maintain training stability and avoid loss spikes. Conversely, higher $\beta_2$ values constrain the learning rate and batch size. Such dynamics have also been recently observed between the learning rate and the first-momentum decay $\beta_1$ in~\citet{pagliardini2025ademamix}. Given that all our experiments use a fixed large batch size of roughly 2 million tokens (appropriate for billion-scale training), we systematically tune the learning rate $\eta$ in response to changes in $\beta_1, \beta_2$. We try values of $\beta_1, \beta_2$ based on previous works~\citep{zhang2025critical} and follow the theoretical convergence requirement of~\citet{zhang2022adam} setting $\beta_1 \leq \sqrt{\beta_2}$.

Due to computational constraints, we cannot jointly optimize synchronization periods, data distributions, and decay parameters, and instead adopt a structured two-stage tuning approach:
\begin{enumerate}
    \item \textbf{Stage 1: Tuning $\eta$ for DDP.} Starting from the recommended baseline learning rate ($\eta_0$)
    from~\citet{allal2025smollm2}, we conduct a grid search as outlined by~\citet{charles2025communication}:
    $\{\ldots, \sqrt{2}^{-2}\eta_0, \sqrt{2}^{-1}\eta_0, \eta_0, \sqrt{2}\eta_0, \sqrt{2}^2\eta_0, \ldots\}$
    We expand this search until perplexity stops improving, identifying an optimal learning rate $\eta^*_{\text{DDP}}$ for each $(\beta_1, \beta_2)$ configuration.

    \item \textbf{Stage 2: Tuning $\eta$ for Local Adam.} We then repeat this procedure for Local Adam,
    using $\eta^*_{\text{DDP}}$ as the new baseline. To balance generalizability and computational cost, we set
    the synchronization period to an intermediate value of $K = 64$, between high-frequency ($K = 16$) and low-frequency ($K = 256$) scenarios.
\end{enumerate}

Additionally, following~\citet{zhang2025critical}, we omit weight decay (set to zero) to simplify the hyperparameter tuning process, as it directly affects only model parameters, not optimizer states.

For experiments using Nesterov, we follow the hyperparameter sweeping procedure of~\citet{charles2025communication}, starting with a server learning rate of $1.0$ and a momentum of $0.9$ and only lowering it if it fails to converge.

\subsubsection{Optimizers' Hyperparameter Configurations}
\label{sec:hp_configs}

\begin{table}[h]
\centering
\caption{Optimal learning rates $\eta^*$ for $\beta_1, \beta_2$ configurations of ADOPT/Adam. The hyperparameter sweep procedure (see Section~\ref{sec:sweep}) involves incrementally adjusting the learning rate by factors of $\sqrt{2}$ around the initial value from~\citet{allal2025smollm2} until performance stops improving.}
\label{tab:lr_sweep}
\begin{tabular}{llll}
\toprule
Optimizer & $\beta_1$ & $\beta_2$ & $\eta^*$ \\
\midrule
\multirow{4}{*}{ADOPT} & 0.9   & 0.9999 & 0.0021 \\
                        & \textbf{0.95}  & \textbf{0.9999} & \textbf{0.0021} \\
                        & 0.99  & 0.9999 & 0.0014 \\
                        & 0.995 & 0.9999 & 0.0007 \\
\midrule
\multirow{5}{*}{Adam}  & 0.9   & 0.95   & 0.0042 \\
                        & \textbf{0.95}  & \textbf{0.95}   & \textbf{0.003}  \\
                        & 0.9   & 0.99   & 0.003  \\
                        & 0.95  & 0.99   & 0.003  \\
                        & 0.99  & 0.99   & 0.0021 \\
\bottomrule
\end{tabular}
\end{table}

Our hyperparameter sweep (Table~\ref{tab:lr_sweep}) indicates that the optimal learning rate $\eta^*$ under the warmup-stable-decay scheduler~\citep{hagele2024scaling} strongly depends on both optimizer type and the chosen $\beta_1, \beta_2$ values. For Adam, optimal learning rates and second-momentum decay ($\beta_2$) align closely with recommendations from~\citet{allal2025smollm2}, though a slightly higher first-momentum decay ($\beta_1$) consistently performs better, in agreement with prior findings~\citep{zhang2025critical}. For ADOPT (default $\beta_2$), we observe a lower optimal learning rate compared to Adam, but similar best-performing $\beta_1$ values. We also find that the optimal learning rate does not differ between DDP and Local Adam for given $\beta_1, \beta_2$ when $K = 64$ and using a $\sqrt{2}$ sweep, higher learning rates either do not provide a benefit or diverge while lower learning rates are only necessary when pushing $K$ far closer to the complete training duration.

We find that increasing $\beta_1$ for ADOPT, and $\beta_1, \beta_2$ for Adam, leads to rapid performance degradation, particularly at or above $0.99$. Since the half-life at $\beta = 0.99$ ($\tau_{0.5} \approx 69$) is not sufficiently longer than at $\beta = 0.95$ ($\tau_{0.5} \approx 13.5$) to justify the observed performance drop, we select $\beta_1 = 0.95$ for all experiments, along with the default $\beta_2$ for ADOPT and $\beta_2 = 0.95$ for Adam.

\paragraph{Takeaway.} Increasing an optimizer state's $\beta$ significantly affects performance. Since linear increases in $\beta$ cause only logarithmic changes in half-life $\tau_{0.5}$, raising $\beta$ beyond the optimal value degrades performance without substantially improving the achievable synchronization frequency (Section~\ref{sec:rq2}).

%=============================================================================
\section{Complementary Results to Sections 2.1 and 5}
\label{sec:complementary}
%=============================================================================

We now provide additional results supplementing those presented in the main text. Specifically:

\begin{enumerate}
\item \textbf{Section B.1} complements Section~2.1 by including results on the heterogeneous data distribution described in Section~4.1. This highlights \texttt{DES-LOC}'s robustness under imperfect sampling or strongly Non-IID federated scenarios (see Kairouz et al., 2021, Sec 3.1).
\item \textbf{Section B.2.1} complements Fig.~3 by showing the separate impact of varying synchronization frequencies for parameters and the second momentum when the base frequency is $K_b=16$. It supports our claim that parameters and second momentum exhibit similar behavior across different synchronization regimes, unlike the first momentum.
\item \textbf{Section B.2.2} extends Fig.~3 by evaluating \texttt{DES-LOC-Adam}. We confirm that the parameter synchronization frequency is the most important, as predicted by our theory. In contrast, the momenta sync frequency is far less impactful, especially for low parameter sync frequencies.
\item \textbf{Section B.3.2} complements Fig.~4 by showing \texttt{DES-LOC-ADOPT}'s perplexity against baseline methods on heterogeneous data (as defined in Section~4.1). This validates our claim from Contribution~2 regarding \texttt{DES-LOC}'s effectiveness on heterogeneous datasets.
\item \textbf{Section B.3.3} presents an ablation study examining alternative low-communication configurations of \texttt{DES-LOC}, justifying our choice of $K_u=3K_x, K_v=6K_x$ used in Fig.~4.
\item \textbf{Section B.3.4} repeats the baseline comparison from Fig.~4 for \texttt{DES-LOC-Adam}, demonstrating that \texttt{DES-LOC} achieves similar communication reductions and performance when using Adam instead of ADOPT.
\item \textbf{Section B.4} provides additional metrics illustrating training instabilities for the \texttt{FAVG+OPT} baseline, including rapidly growing parameter norms, supporting observations in Fig.~5.b.
\item \textbf{Section B.5} presents the full 1B-scale comparison of \texttt{DES-LOC} Nesterov against DiLoCo, demonstrating the benefits of periodic optimizer-state synchronization under a Nesterov outer optimizer.
\item \textbf{Section B.6} benchmarks \texttt{DES-LOC} under extremely low bandwidth ($10$ Gbit/s), showing $9.42\times$ wall-clock speedups over DDP.
\item \textbf{Section B.7} provides full experimental details for the Muon experiments presented in Section~5.6.
\item \textbf{Section B.8} evaluates \texttt{DES-LOC} on the Flux vision model, demonstrating generality beyond LLMs.
\item \textbf{Section B.9} measures training throughput at the $7$B parameter scale on B200 GPUs.
\end{enumerate}

%-----------------------------------------------------------------------------
\subsection{Toy Problem on Non-IID Data (See Section~2.1)}
\label{sec:b1_toy_noniid}
%-----------------------------------------------------------------------------

\textbf{Toy Example Non-IID:} Fig.~8 simulates the scenario from Section~3, where each worker $m$ optimizes a distinct loss $f_m$ on heterogeneous data. Both \texttt{DES-LOC} and Local Adam show more stable convergence and get closer to the optimum than heuristic baselines.

% Figure 8: Toy problem Non-IID (PDF page 21 top)
\begin{figure}[ht]
\centering
% [placeholder for toy problem Non-IID figure]
\caption{We present a toy problem in a Non-IID setting, where \texttt{DES-LOC} (with synchronization periods $K_x=192, K_u=192, K_v=692$) and Local Adam (with $K=K_x$) converge to a superior solution compared to methods that keep optimizer states local \citep{Douillard2023,Sani2025} or periodically reset them \citep{Sani2024,Iacob2025}. Like the IID scenario, resetting optimizer states prevents convergence due to repeated oscillations caused by reinitializations. Additionally, as seen in panel~(a) between rounds $15$ and $40$, methods keeping optimizer states local suffer from larger oscillations further away from the optimum. The function optimized is $f(x_1,x_2)=(1-x_1)^2+100(x_2-x_1^2)^2$, and we simulate $M=256$ workers, each adding Gaussian noise with worker-specific standard deviation $\sigma^m \sim \mathcal{N}(0,3)$.}
\label{fig:8}
\end{figure}

%-----------------------------------------------------------------------------
\subsection{RQ2: Independent Sync Frequencies}
\label{sec:b2_sync_freq}
%-----------------------------------------------------------------------------

This section provides supplementary results for \textbf{RQ2}, complementing Section~5.2. Section~B.2.1 shows that perplexity has similar sensitivity to the first and second momentum synchronization frequencies at both high and low base synchronization frequencies. Additionally, Section~B.2.2 repeats the comparison from Fig.~3 for \texttt{DES-LOC-Adam}, revealing similar trends regarding the importance of the parameters, with a reduced importance for the momenta due to lower $\beta_2$.

\subsubsection{Parameter and Second Momentum At $K_b=16$ (See Fig.~3.a, Fig.~3.b)}
\label{sec:b21_param_second}

Figure~9 examines the effects of independently varying synchronization periods ($K_x, K_v$) for parameters and second momentum under \texttt{DES-LOC-ADOPT} in the high-frequency regime ($K_b=16$), chosen based on the first momentum's half-life ($\tau_{0.5}\approx 13.5$). Similar to the low-frequency results in Fig.~3.a, parameter synchronization frequency ($K_x$) strongly influences perplexity, while the second momentum ($K_v$) has minimal impact due to its long half-life. This contrasts with the first momentum, whose half-life closely matches the high-frequency period.

% Figure 9: DES-LOC ADOPT perplexity at Kb=16 (PDF page 21 bottom)
\begin{figure}[ht]
\centering
% [placeholder for Fig 9]
\caption{Model perplexity for \texttt{DES-LOC} (ADOPT, $\beta_1=0.95, \beta_2=0.9999$), independently varying synchronization periods at a high baseline frequency ($K_b=16$). Similar to Fig.~3, parameter synchronization~(a) is critical, with performance sharply degrading at higher periods, while second-momentum synchronization~(b) has minimal impact due to its large half-life ($\tau_{0.5}(\beta_2)\gg K_b$).}
\label{fig:9}
\end{figure}

\textbf{Takeaway:} In high-frequency synchronization regimes, the importance of parameters and the second momentum remains similar to the low-frequency regime shown in Section~5.2.

\subsubsection{Adam Results (See Fig.~3)}
\label{sec:b22_adam}

Fig.~10 show the rate of change results for Adam momenta at various $\beta$ values.

Figure~11 provides complementary results to Figs.~3 and~9 using \texttt{DES-LOC-Adam} with $\beta_1=\beta_2=0.95$. Unlike ADOPT, the relatively low $\beta$ result in both the first and second momentum quickly adapting to the local gradients, reducing the impact of their sync frequency.

% Figure 10: Rate of change for Adam momenta (PDF page 22)
\begin{figure}[ht]
\centering
% [placeholder for Fig 10]
\caption{Relative rates of change for first and second momenta across rounds using standard Local Adam ($K=64$). Increasing $\beta_1$ substantially reduces the rate of change of the first momentum, while increasing either $\beta_1, \beta_2$ decreases the rate of change of the second.}
\label{fig:10}
\end{figure}

% Figure 11: DES-LOC-Adam perplexity varying sync periods (PDF page 22)
\begin{figure}[ht]
\centering
% [placeholder for Fig 11]
\caption{Model perplexity for \texttt{DES-LOC-Adam} ($\beta_1=\beta_2=0.95$) when independently varying sync periods ($K_x, K_u, K_v$) while fixing others at baseline $K_b$. Parameter synchronization~(a,b) influences performance in both high ($K_b=16$) and low ($K_b=256$) frequency regimes. Momenta synchronization minimally impacts perplexity due to both states' high adaptivity (low $\beta$), with potentially minor effects during the early stages of training in high-frequency regimes~(c,e).}
\label{fig:11}
\end{figure}

\textbf{Takeaway:} For \texttt{DES-LOC-Adam}, parameter synchronization remains critical, consistent with theory. However, due to reduced $\beta_2$, momenta synchronization is less impactful since both the numerator and denominator of Adam updates are driven by local worker gradients after a few initial steps.

%-----------------------------------------------------------------------------
\subsection{RQ3: Communication Reduction And Baseline Comparisons}
\label{sec:b3_comm_reduction}
%-----------------------------------------------------------------------------

This section provides supplementary results for \textbf{RQ3}, complementing Section~5.3. Section~B.3.2 shows the perplexity of different configurations providing a $2\times$ communication reduction over Local Adam. Additionally, Section~B.3.4 repeats the comparison against baselines from Section~5.3 for \texttt{DES-LOC-Adam}, showing similar communication reductions relative to Local Adam.

\subsubsection{Stepwise plots for Baseline Comparison}
\label{sec:b31_stepwise}

Figures~12 and~13 show stepwise plots for wall-clock results in the main text, they are the counterparts to Figs.~4 and~5.

% Figure 12 and 13: Stepwise plots (PDF page 23)
% [placeholder for Figs 12-13 -- stepwise counterparts to Figs 4,5]

\subsubsection{DES-LOC on Heterogeneous Data (See Contribution~2)}
\label{sec:b32_heterogeneous}

Figure~14 evaluates the robustness of \texttt{DES-LOC} against baselines under heterogeneous (Non-IID) data distributions as described in Section~4.1. We set synchronization periods to $K_x=K$, $K_u=3K_x$, and $K_v=6K_x$ to achieve a targeted $2\times$ communication reduction over Local Adam.

% Figure 14: Non-IID perplexity comparison (PDF page 24 top)
\begin{figure}[ht]
\centering
% [placeholder for Fig 14]
\caption{Comparison of perplexity under Non-IID conditions for \texttt{DES-LOC}, Local Adam ($K_x=K_u=K_v$), and heuristic baselines (defined in Section~4.1) at high~(a) and low~(b) synchronization frequencies. Due to higher cross-worker variance caused by heterogeneous data, parameters require slightly more frequent synchronization in the low-frequency regime ($K_x=128<256$). Experiments are limited to $T=1536$ steps ($\sim$compute-optimal) for computational feasibility.}
\label{fig:14}
\end{figure}

\subsubsection{DES-LOC Low Communication Configurations Ablation (See Fig.~4)}
\label{sec:b33_ablation}

Figure~15 explores alternative synchronization configurations enabling \texttt{DES-LOC} to achieve improved communication efficiency over Local Adam. Motivated by theoretical insights (Sections~2 and~3) and empirical evidence (Sections~5.1 and~5.2), we only consider settings where parameter synchronization is most frequent ($K_x \leq \min(K_u, K_v)$). This constraint follows from experiments in Section~5.2, which show that infrequent parameter synchronization significantly degrades perplexity, while momentum synchronization frequency has a smaller impact. For a fixed $2\times$ communication reduction over Local Adam, our findings confirm that synchronizing the first momentum more frequently than the second aligns with their respective half-lives and maintains performance close to Local Adam.

% Figure 15: DES-LOC 2x communication configs ablation (PDF page 24 bottom)
\begin{figure}[ht]
\centering
% [placeholder for Fig 15]
\caption{Configurations of \texttt{DES-LOC} targeting $2\times$ lower communication than Local Adam ($K_x=K_u=K_v$), setting $K_u, K_v$ as multiples of $K_x$. In both high~(a) and low-frequency~(b) regimes, performance depends on how communication is split between momenta for $\beta_1 \ll \beta_2$. Syncing the first momentum less often ($K_u=6K_x, K_v=3K_x$) degrades performance, wasting communication on the slow second momentum. Conversely, syncing it frequently ($K_u=3K_x, K_v=6K_x$) yields performance comparable to Local Adam. Setting $K_u=K_v=4K_x$ produces intermediate results.}
\label{fig:15}
\end{figure}

\subsubsection{Adam Results (See Fig.~4)}
\label{sec:b34_adam_fig4}

We now present results for \texttt{DES-LOC-Adam} with $\beta_1=\beta_2=0.95$. \texttt{DES-LOC-Adam} achieves similar communication reductions over Local Adam and DDP as ADOPT. However, due to the lower $\beta_2$, the second-momentum half-life ($\tau_{0.5}(0.95)\approx 13.5$) is significantly shorter than for ADOPT ($\tau_{0.5}(0.9999)\approx 6931$). Figure~16 shows that with both momenta evolving at similar rates, the benefit of selecting $K_u < K_v$ diminishes. For consistency and due to meaningful empirical differences in rates of change (Section~5.1), we keep $K_u=3\times K_x$ and $K_v=6\times K_x$ in subsequent comparisons.

Figure~17 shows \texttt{DES-LOC-Adam} achieves a $2\times$ communication reduction over the prior state-of-the-art Local Adam \citep{Cheng2025} without significant perplexity degradation. Due to the much faster evolution of the optimizer states using Adam compared to ADOPT, local worker gradients drive the optimization reducing the benefit of allocating more of the communication budget to the first momentum.

% Figure 16: DES-LOC-Adam 2x communication configs (PDF page 25)
\begin{figure}[ht]
\centering
% [placeholder for Fig 16]
\caption{Configurations of \texttt{DES-LOC} targeting $2\times$ lower communication than Local Adam ($K_x=K_u=K_v$), using Adam ($\beta_1=\beta_2=0.95$). In contrast to \texttt{DES-LOC-ADOPT} (where $\beta_1 \ll \beta_2$ yields an advantage for $K_u < K_v$ as shown in Fig.~15), the similar half-lives in Adam make perplexity insensitive to how communication is split between momenta for high~(a) and low-frequencies~(b).}
\label{fig:16}
\end{figure}

% Figure 17: DES-LOC-Adam baseline comparison (PDF page 25)
\begin{figure}[ht]
\centering
% [placeholder for Fig 17]
\caption{\texttt{DES-LOC-Adam} achieves a $2\times$ communication reduction over the prior state-of-the-art Local Adam without significant perplexity degradation.}
\label{fig:17}
\end{figure}

\textbf{Takeaway:} \texttt{DES-LOC-Adam} achieves a similar $2\times$ communication reduction over Local Adam as \texttt{DES-LOC-ADOPT} by exploiting the reduced importance of optimizer-state synchronization relative to parameters. However, due to the smaller $\beta_2$ in Adam, there is limited benefit from assigning different synchronization frequencies to the first and second momenta compared to ADOPT.

%-----------------------------------------------------------------------------
\subsection{RQ4: Additional Metrics and Training Instabilities of FAVG+OPT (See Fig.~5.b)}
\label{sec:b4_favgopt}
%-----------------------------------------------------------------------------

Figure~18 complements Fig.~5.b by showing parameter and update norms for \texttt{DES-LOC} and baseline methods when training billion-scale models. Both \texttt{DES-LOC} and Local Adam regularize updates by synchronizing optimizer states, effectively reducing update norms due to averaging across workers (triangle inequality). In contrast, the heuristic baseline \citep{Sani2025} experiences large updates, leading to uncontrolled parameter growth, increased activations (Fig.~5.b), and degraded performance on downstream ICL tasks (Table~1) relative to its perplexity (Fig.~5.a).

% Figure 18: Update and parameter norms comparison (PDF page 26)
\begin{figure}[ht]
\centering
% [placeholder for Fig 18]
\caption{Comparison of update~(a) and parameter norms~(b) for billion-scale models trained with \texttt{DES-LOC} ($K_x=256, K_u=768, K_v=1536$), Local Adam ($K=256$), DDP, and Federated Averaging with persistent optimizer states (\texttt{FAVG+OPT}). Frequent synchronization in Local Adam and DDP consistently reduces update and parameter norms. Similarly, \texttt{DES-LOC} achieves comparable reductions at intervals corresponding to multiples of $\texttt{lcm}(K_x, K_u, K_v)$, with smaller intermediate drops. Conversely, \texttt{FAVG+OPT}, which does not synchronize optimizer states, experiences persistently larger and noisier updates, becoming vulnerable to spikes (notably before step $5000$). This leads to uncontrolled parameter growth~(b).}
\label{fig:18}
\end{figure}

% Table 5: Throughput comparison (PDF page 26, nougat inline)
\begin{table}[ht]
\centering
\caption{Throughput (batches/sec) for 1B--13B models to reach $2\times$ compute-optimal tokens at high ($K=16$) and low ($K=256$) frequencies with a 2M token batch size. All local methods achieve significant throughput gains over the DDP baseline. At high frequency ($K=16$), \texttt{DES-LOC} boosts throughput by over $1.7\times$ on the $13$B model. At low frequency ($K=256$), this advantage grows to over $2.1\times$ versus DDP.}
\label{tab:5}
\begin{tabular}{l c c c c c c}
\hline \hline
 & \multicolumn{2}{c}{\textbf{1B Model}} & \multicolumn{2}{c}{\textbf{7B Model}} & \multicolumn{2}{c}{\textbf{13B Model}} \\
\cline{2-7}
\textbf{Method} & $K_x\!=\!16$ & $K_x\!=\!256$ & $K_x\!=\!16$ & $K_x\!=\!256$ & $K_x\!=\!16$ & $K_x\!=\!256$ \\
\hline
DDP (Baseline) & $171.9 \pm 0.94$ & $171.9 \pm 0.9$ & $25.1 \pm 0.10$ & $25.1 \pm 0.10$ & $11.1 \pm 0.03$ & $11.11 \pm 0.03$ \\
\hline
FAVG+OPT & $304.9 \pm 2.51$ & $385.4 \pm 3.7$ & $34.0 \pm 0.11$ & $40.4 \pm 0.15$ & $19.4 \pm 0.11$ & $23.5 \pm 0.15$ \\
Local Adam & $253.1 \pm 1.59$ & $380.1 \pm 3.6$ & $30.9 \pm 0.09$ & $40.2 \pm 0.15$ & $16.7 \pm 0.07$ & $23.3 \pm 0.14$ \\
DES-LOC ($K_u,K_v\!=\!3K_x,6K_x$) & $299.2 \pm 2.39$ & $384.1 \pm 3.7$ & $33.7 \pm 0.11$ & $40.4 \pm 0.15$ & $19.0 \pm 0.10$ & $23.5 \pm 0.15$ \\
\hline \hline
\end{tabular}
\end{table}

\textbf{Takeaway:} Unlike heuristic methods, which maintain purely local optimizer states leading to unstable, noisy updates, \texttt{DES-LOC} provides stable regularization similar to Local Adam and DDP by periodically synchronizing parameters and momenta, reducing training instabilities.

% B.5-B.9: Deferred to Claude-9 (M201-M225)
% B.5: DES-LOC Nesterov vs DiLoCo at 1B → DES-LOC competitive
% B.6: Very low bandwidth (10 Gbit/s) → DES-LOC still beneficial
% B.7: Muon experimental details → Kx=32, Ku=96
% B.8: Flux vision model → DES-LOC generalizes beyond LLMs
% B.9: 7B throughput on B200 → near-linear scaling

%=============================================================================
\section{Further Algorithmic Details of DES-LOC}
\label{sec:appendix_c}
%=============================================================================

\subsection{Extension to FedOpt}
\label{sec:fedopt}

Although~\citet{cheng2025convergence} show provable convergence for adaptive inner optimizers in a
federated optimization framework, their result rests on the assumption that after a period of local
work, the new global model is created by averaging the local client models. In relation to the larger
FedOpt literature~\citep{reddi2021adaptive}, the scheme chosen by~\citet{cheng2025convergence} resembles
that of FedAvg, or where the server optimizer is SGD with the outer learning rate set to one.

We argue that indeed DES-LOC's principles can be effectively applied to any FedOpt method and
not just FedAvg. While using an alternate server optimizer does not have proven convergence
guarantees as yet, we show in Algorithm~\ref{alg:desloc} that the choice of the ServerOpt is not constrained
from a practical point-of-view. However, the improvements that DES-LOC provide are related to
the local optimization procedure, which is orthogonal to the outer optimizer choice.

\subsection{Deterministic Optimizer-specific Variants of Algorithm 1}
\label{sec:algorithm_variants}

\begin{algorithm}[h]
\caption{DES-LOC-Adam}
\label{alg:desloc_adam}
\begin{algorithmic}[1]
\REQUIRE Model tensors, Hyper-parameters
\STATE $x_0, u_{-1}, v_{-1} \in \mathbb{R}^d$ \COMMENT{initial parameter vector, seeds for first and second moments}
\STATE $\{\eta_t\}_{t=0}^{T-1} \subset \mathbb{R}_{>0}$ \COMMENT{step-size schedule}
\STATE $\beta_1, \beta_2 \in [0,1)$ \COMMENT{Adam decay factors}
\STATE $\rho, \lambda \in \mathbb{R}_{>0}$ \COMMENT{gradient clipping term, $\ell_2$ stability term}
\STATE $T, M \in \mathbb{N}_+$ \COMMENT{total iterations, number of workers}
\STATE $K_x, K_u, K_v \in \mathbb{N}_+$ \COMMENT{sync periods for parameters, first and second moments}
\ENSURE $x_T, u_{T-1}, v_{T-1}$
\STATE for each worker $m$: $x^m_0 = x_0$, $u^m_{-1} = v^m_{-1} = 0$
\FOR{$t = 0, \ldots, T-1$}
    \FOR{all workers $m = 0, \ldots, M-1$ \textbf{in parallel}}
        \STATE $g^m_t \gets \nabla F(x^m_t; \xi^m_t)$
        \STATE $\hat{g}^m_t \gets \text{clip}(g^m_t, \rho)$
        \IF{$t \bmod K_u = 0$}
            \STATE $u^m_t \gets \beta_1 \mathbb{E}_m[u^m_{t-1}] + (1 - \beta_1)\hat{g}^m_t$
        \ELSE
            \STATE $u^m_t \gets \beta_1 u^m_{t-1} + (1 - \beta_1)\hat{g}^m_t$
        \ENDIF
        \IF{$t \bmod K_v = 0$}
            \STATE $v^m_t \gets \beta_2 \mathbb{E}_m[v^m_{t-1}] + (1 - \beta_2)(\hat{g}^m_t \odot \hat{g}^m_t)$
        \ELSE
            \STATE $v^m_t \gets \beta_2 v^m_{t-1} + (1 - \beta_2)(\hat{g}^m_t \odot \hat{g}^m_t)$
        \ENDIF
        \STATE $m^m_t \gets \beta_1 m^m_{t-1} + (1 - \beta_1) \frac{\hat{g}^m_t}{\max\{\sqrt{v^m_{t-1}}, \epsilon\}}$ \COMMENT{ADOPT update}
        \STATE $d^m_t \gets \eta_t m^m_t$
        \IF{$t \bmod K_x = 0$}
            \STATE $x^m_{t+1} \gets \mathbb{E}_m[x^m_t] - d^m_t$
        \ELSE
            \STATE $x^m_{t+1} \gets x^m_t - d^m_t$
        \ENDIF
    \ENDFOR
\ENDFOR
\end{algorithmic}
\end{algorithm}

%=============================================================================
% Appendix D: Convergence Analysis of DES-LOC-SGDM
% NOTE: Full 8-page proof deferred to Claude instances 12-13 (M276-M325)
%=============================================================================
\section{Convergence Analysis of DES-LOC-SGDM (in expectation bounds)}
\label{sec:sgdm_proof}
% Placeholder — full proof spans pages 32-43 of the PDF

%=============================================================================
% Appendix E: High-probability bounds for DES-LOC-Adam
% NOTE: Full proof deferred to Claude instance 14 (M326-M350)
%=============================================================================
\section{Convergence Analysis of DES-LOC-Adam (high-probability bounds)}
\label{sec:high_prob}
% Placeholder — full proof spans pages 44-46 of the PDF

%=============================================================================
\section{Derivation of Eqs.\ (1) and (2): Maximum Momentum Change With Clipping}
\label{sec:drift_derivation}
%=============================================================================

\textbf{Step 2: Bound on $\|u_{t+K} - u_t\|_\infty$.} Now we bound the change in the momentum over $K$ steps
explicitly. Unrolling the recursion, we have:
\begin{equation}
u_{t+K} = \beta_1^K u_t + (1 - \beta_1) \sum_{k=0}^{K-1} \beta_1^k g_{t+K-k}.
\label{eq:unroll_u}
\end{equation}

Subtracting $u_t$ from both sides, we obtain:
\begin{equation}
u_{t+K} - u_t = (\beta_1^K - 1)u_t + (1 - \beta_1) \sum_{k=0}^{K-1} \beta_1^k g_{t+K-k}.
\label{eq:diff_u}
\end{equation}

Applying the triangle inequality gives:
\begin{equation}
\|u_{t+K} - u_t\|_\infty \leq |1 - \beta_1^K| \|u_t\|_\infty + (1 - \beta_1) \sum_{k=0}^{K-1} \beta_1^k \|g_{t+K-k}\|_\infty.
\label{eq:triangle_u}
\end{equation}

Using $\|g_{t,i}\|_\infty \leq \rho$ (from coordinate-wise clipping) and $\|u_t\|_\infty \leq \rho$ (by induction, since $u$ is an exponential moving average of clipped gradients), the geometric series evaluates to:
\begin{equation}
\|u_{t+K} - u_t\|_\infty \leq (1 - \beta_1^K)\rho + (1 - \beta_1^K)\rho = 2\rho(1 - \beta_1^K).
\end{equation}

The identical argument applies to the second moment $v_t$ with $\rho$ replaced by $\rho^2$ (since $(g_t \odot g_t)_i \leq \rho^2$), yielding:
\begin{equation}
\|v_{t+K} - v_t\|_\infty \leq 2\rho^2(1 - \beta_2^K).
\end{equation}

%=============================================================================
\section{Wall-Clock Time Modeling}
\label{sec:wallclock_model}
%=============================================================================

\subsection{Estimating Total Wall-Clock Time}

We model the total wall-clock time $\mathcal{T}$ for training as:
\begin{equation}
\mathcal{T} = T \cdot t_{\text{step}} + R_x \cdot t_{\text{comm}}(x) + \sum_{j=1}^{N} R_j \cdot t_{\text{comm}}(s^j),
\end{equation}
where $T$ is the total number of steps, $t_{\text{step}}$ is the per-step computation time, $R_x = T/K_x$ is the number of parameter sync rounds, $R_j = T/K_j$ is the number of sync rounds for state $j$, and $t_{\text{comm}}(\cdot)$ is the communication time for a given tensor.

For Ring-AllReduce with model size $|x|$ parameters, bandwidth $B_w$ (bytes/sec), and $M$ workers:
\begin{equation}
t_{\text{comm}}(x) = \frac{2(M-1)}{M} \cdot \frac{|x| \cdot \text{bytes\_per\_param}}{B_w}.
\end{equation}

DES-LOC's communication cost relative to DDP is:
\begin{equation}
\frac{\mathcal{T}_{\text{DES-LOC}}}{\mathcal{T}_{\text{DDP}}} = \frac{T \cdot t_{\text{step}} + \frac{T}{K_x} t_{\text{comm}}(x) + \frac{T}{K_u} t_{\text{comm}}(u) + \frac{T}{K_v} t_{\text{comm}}(v)}{T \cdot (t_{\text{step}} + t_{\text{comm}}(x) + t_{\text{comm}}(u) + t_{\text{comm}}(v))}.
\end{equation}

With $K_x = 256$, $K_u = 3K_x = 768$, $K_v = 6K_x = 1536$, DES-LOC communicates $\approx 170\times$ less than DDP.

\begin{table}[h]
\centering
\caption{Throughput (tokens/sec/GPU) at different scales. DES-LOC matches FAVG+OPT throughput.}
\label{tab:throughput}
\begin{tabular}{l|cc|cc}
\toprule
& \multicolumn{2}{c|}{1B Model} & \multicolumn{2}{c}{7B Model} \\
$K_x$ & 16 & 256 & 16 & 256 \\
\midrule
DDP          & 60k & 60k & 8.5k & 8.5k \\
FAVG+OPT     & 106k & 135k & 11.5k & 13.7k \\
Local Adam   & 89k & 132k & 10.4k & 13.6k \\
DES-LOC      & 105k & 134k & 11.4k & 13.7k \\
\bottomrule
\end{tabular}
\end{table}

%=============================================================================
\section{Checkpointing vs.\ Periodic State Synchronization}
\label{sec:checkpoint}
%=============================================================================

A natural question is whether periodic checkpointing can replace DES-LOC's explicit state synchronization for handling worker failures. We argue it cannot, for two reasons:

\begin{enumerate}
\item \textbf{State staleness.} When a new worker joins from a checkpoint, its optimizer states ($u, v$) reflect the global state at checkpoint time. If the checkpoint is old, these states are stale and can cause training instability. DES-LOC's periodic synchronization ensures all workers' states stay within bounded drift (Eqs.~\ref{eq:drift_u},~\ref{eq:drift_v}).

\item \textbf{No convergence guarantees.} Ad-hoc checkpoint averaging provides no theoretical guarantees about the quality of the resulting optimizer states. DES-LOC's provable convergence (Theorem~\ref{thm:sgdm}) holds precisely because synchronization is periodic and bounded.
\end{enumerate}

%=============================================================================
\section{Choosing Synchronization Frequencies}
\label{sec:choosing_freq}
%=============================================================================

For a state with a decay rate $\beta$ approaching 1, the half-life can be calculated as:
\begin{equation}
t_{1/2} \approx \frac{\ln 2}{1 - \beta}.
\end{equation}

Based on this, we propose the following two-step methodology:

\paragraph{Parameters First ($K_x$).} The synchronization of model parameters is paramount to training
quality. The period $K_x$ should be chosen to match the end-of-training quality of fully-synchronous DDP at a target step budget while still materially reducing communication. In
practice, starting points like $K_x = 16$ or $K_x = 32$ are effective.

\paragraph{Momentum by Half-Life ($K_u, K_v$).} For any optimizer momentum state with a decay rate
$\beta$, its synchronization period $K$ should be set near its calculated half-life, i.e., $K \approx t_{1/2}$.
For common optimizers like Adam or ADOPT with well-tuned decay rates $\beta_1$ and $\beta_2$, this
simplifies to:
\begin{equation}
K_u \approx \frac{\ln 2}{1 - \beta_1}, \quad K_v \approx \frac{\ln 2}{1 - \beta_2}.
\end{equation}

Following this heuristic can yield a minimum $5\times$ reduction in communication cost over DDP for
$K_x \geq 16$, significantly decreasing wall-clock time while achieving convergence speed and final
model quality comparable to DDP.

%=============================================================================
\section{Critical Batch Size and Regime Positioning}
\label{sec:critical_batch}
%=============================================================================

To formally contextualize the regimes where DES-LOC is most beneficial, we begin with the statistical
properties of gradient estimation. Let the true gradient over the full data distribution for a loss function
$L: \mathbb{R}^d \to \mathbb{R}$ be $G(\theta) = \nabla L(\theta)$. In practice, a mini-batch of size $B$ provides an estimate, $G_{\text{est}}(\theta)$.
The variance of this estimator scales inversely with the batch size:
\begin{equation}
\text{Var}[G_{\text{est}}(\theta)] = \frac{1}{B}\Sigma(\theta),
\end{equation}
where $\Sigma(\theta)$ is the per-sample gradient covariance. DES-LOC is most beneficial in the large-batch regime where communication costs dominate compute costs.

%=============================================================================
\section{Extended Related Work}
\label{sec:extended_related}
%=============================================================================

\paragraph{Gradient compression and sparsification.}
Methods like Deep Gradient Compression~\citep{lin2018deep}, TopK sparsification~\citep{alistarh2018convergence}, and CocktailSGD~\citep{wang2023cocktailsgd} reduce communication by compressing or sparsifying gradients. These are orthogonal to DES-LOC and can be combined for further savings.

\paragraph{Model parallelism.}
Megatron-LM~\citep{shoeybi2019megatron}, ZeRO~\citep{rajbhandari2020zero}, and FSDP~\citep{zhao2023pytorch} partition model parameters across devices. DES-LOC focuses on data parallelism and is complementary to model-parallel approaches.

\paragraph{Personalized and heterogeneous federated learning.}
Methods like Ditto~\citep{li2021ditto}, FedProx~\citep{li2020federated}, and Per-FedAvg~\citep{fallah2020personalized} address non-IID data in federated settings. DES-LOC inherits robustness to heterogeneity from its theoretical framework (Assumption~\ref{asm:heterogeneity}).

%=============================================================================
\section{LLM Usage Declaration}
\label{sec:llm_usage}
%=============================================================================

Large Language Models were used during the preparation of this work for editing and improving the
text quality of the manuscript. No LLM was used for generating scientific content, proofs, or
experimental results. All technical contributions are the original work of the authors.

%=============================================================================
\section{Limitations}
\label{sec:limitations}
%=============================================================================

\begin{enumerate}
\item \textbf{Communication overlap.} Our current implementation uses a ``stop-the-world'' synchronization approach. Overlapping communication with computation~\citep{douillard2025streaming} could further improve wall-clock time.

\item \textbf{Heterogeneous hardware.} We assume homogeneous worker hardware. Extending DES-LOC to heterogeneous clusters with different compute speeds is left for future work.

\item \textbf{Adaptive synchronization.} Our synchronization periods are fixed throughout training. Adapting them based on training dynamics (e.g., increasing $K_x$ as training progresses) could yield further improvements.

\item \textbf{Scale of experiments.} While we validate up to 1.7B parameters and project to 13B, experiments at truly large scale (100B+) remain future work due to resource constraints.
\end{enumerate}
